% Appendices for the divide-and-conquer half. Needs \usepackage{listings}.

\section{Divide-and-Conquer Code}
\label{app:dc-code}

The algorithm (\texttt{quadtree.py}) and the measurements behind
Figs.~\ref{fig:dc-scaling}--\ref{fig:dc-depth} and Table~\ref{tab:dc-requests}
(\texttt{experiments.py}). The plotting script, the unit tests, the Rinc\'on
data file and the raw measurements are in the project
repository.\footnote{\url{https://github.com/IlanDanial/P1Alg}}

\lstinputlisting[language=Python,caption={quadtree.py}]{\dcdir quadtree.py}
\lstinputlisting[language=Python,caption={experiments.py}]{\dcdir experiments.py}

\section{Machine-Checked Proof}
\label{app:dc-lean}

Lemma~\ref{lem:partition} (\texttt{quadrants\_cover},
\texttt{quadrants\_disjoint}), Theorem~\ref{thm:correct}
(\texttt{collect\_correct}), Lemma~\ref{lem:height} on the integer grid
(\texttt{collect\_terminates}) and their combination into total correctness
(\texttt{collect\_total}) in Lean~4.35, using only the core library.
\texttt{lake build} in \texttt{divide\_and\_conquer/proof} checks it; the
theorems depend on no axioms beyond Lean's standard three
(\texttt{propext}, \texttt{Classical.choice}, \texttt{Quot.sound}).

\lstinputlisting[language={},caption={proof/Proof/Collect.lean}]{\dcdir proof/Proof/Collect.lean}

\section{LLM Usage: Divide-and-Conquer Part}
\label{app:dc-llm}

\textbf{Tool.} Claude Code (Anthropic's command-line coding agent), model
Claude Opus~5.5 (\texttt{claude-opus-5-5}), October 2026. The agent read the
assignment, read the scraper repository~\cite{danial2023airbnb}, searched the web
for the IEEE template, wrote and ran the code, and drafted this section. All
proofs, numbers, and citations were checked by the authors.

\textbf{Prompts (verbatim).}
\begin{enumerate}
\item ``DO you understand the project I need to do also this is a group project
  I am doing the divide an conquer part''
\item ``My Idea for the problem is a problem I have solved before with divide
  and conquer \url{https://github.com/IlanDanial/airbnb_data} I split the map over and
  over, each time checking the amount of properties in the split and if it was
  under I took those properties and added it to the list since you can only see
  so many properties at a time''
\item ``yes start the implementation, we're using IEEE search online for proper
  template''
\item ``yes write the LaTeX section''
\item ``This is really ugly'' followed by ``My bigger issue is that the
  explanations are very longwinded for no reason and the formatting of the
  page is ugly its hard to read as well''
\item ``I have an issue with some of this read the paper be critical how do
  you think of the formatting'', then ``fix all of them at once'' and ``also
  do we need the full code in the appendix or can we reference a repository''
\item ``would it be worth it to look at the code and refactor it and the
  comments in the code'', then ``go''
\item ``Isnt there a code languge that is built to prove proofs can we use
  that to prove correctness'', then ``I want to go for it also I am sure there
  are online proofs of this exact same thing in the same language so it should
  take less than an hour''
\item ``so your telling me that there was no problem with proving the whole
  thing end to end''
\item ``actually search online for the hard part dont just tell me it does
  not exist''
\item ``do both, go for the full proof, or is there a solution we can do that
  would be easier to prove and solve the problem all the same''
\end{enumerate}

\textbf{Results.}
% TODO: paste the full exported transcript (Claude Code: /export) here or in
% supplementary material; the assignment requires all intermediate results.
(1)~A summary of the assignment and four candidate divide-and-conquer problems.
(2)~An analysis of the scraper's \texttt{subdivide} recursion as an adaptive
quadtree, identifying the termination issue for co-located listings and the
double-counting issue on window borders.
(3)~\texttt{quadtree.py}, \texttt{test\_quadtree.py}, \texttt{experiments.py},
and the figures; an intermediate run stopped when the minimum-width guard was
reached, after which the crowding experiment was limited to
$\delta\ge2^{-28}$, and garbage collection was disabled during timing.
(4)~A first draft of Section~\ref{sec:dc} and these appendices, with five
separate figures in matplotlib's default style.
(5)~The draft was shortened by about a third, the five figures were
restyled to match the paper's typeface and merged into three, and the page
layout was fixed.
(6)~A list of nine formatting defects (an equation label pushed to its own
line, a wrapped pseudocode line, mixed paragraph styles, inconsistent
equation citations, a legend mismatch), all then fixed; the unit tests were
removed from this appendix in favor of the repository link.
(7)~The plotting code was split out of \texttt{experiments.py} into
\texttt{plots.py}, and docstrings and comments were added.
(8)~An assessment of proof assistants (Lean~4, Coq, Isabelle, Dafny) and of
which parts of the proof are feasible to formalize; then a first Lean
formalization of Lemma~\ref{lem:partition} and Theorem~\ref{thm:correct}
only, with \textsc{Collect} given a fuel bound, which compiled on the first
build.
(9)~A clarification that this was partial correctness, not an end-to-end
proof, with three further gaps listed (integer midpoints, the oracle not being
a black box, no duplicate-freeness proof).
(10)~A web search that found the Isabelle and Coq formalizations cited in
Section~\ref{sec:dc-correctness}, none covering separation-based termination;
the assistant's earlier claim that no related formalization existed was
withdrawn.
(11)~The observation that finite-precision coordinates make the problem a
grid problem, on which Lemma~\ref{lem:height} is provable by induction on the
depth; \texttt{collect\_terminates} and \texttt{collect\_total} were then
added, needing one round of fixes to core-library lemma names before
compiling.
